\hypertarget{appendix-7679107d057e}{}
\section*{Rules of Three Five and Zero}
\section{Rules of Three Five and Zero}\label{rules-of-three-five-and-zero}

\textbf{Curriculum depth:} Introduced · \textbf{First seen:} C009 · \textbf{Latest development:} C009

\subsection{Definition and mechanism}\label{definition-and-mechanism}

Three coordinates destruction/copy construction/copy assignment; Five adds both moves. These are semantic design decisions, not requirements to implement copying. Zero composes members with correct ownership so enclosing classes avoid custom transfer/cleanup boilerplate; domain invariants still need design.

\subsection{Prerequisites}\label{prerequisites}

\begin{itemize}
\tightlist
\item
  \hyperref[appendix-42e02a12c694]{Special member generation and traits} --- The actual usable construction expression determines trait-controlled transfer selection.
\item
  \hyperref[appendix-bea6beb5350f]{Assignment and exception guarantees} --- A type's transfer declarations must support its chosen failure transaction.
\item
  \href{}{Ownership and non-owning views (05-Engine-I, appendix-d52e4486f527)} --- The owner/view distinction determines whether copying duplicates resources or descriptors.
\end{itemize}

\subsection{Used by}\label{used-by}

Later sources will extend this node when their guarantees are developed.

\subsection{Recall and next depth}\label{recall-and-next-depth}

Replace the teaching raw byte owner with a vector and identify remaining invariant obligations. Pending: reusable unique GPU/OS owners and error policies. Written coverage does not demonstrate personal mastery.

\subsection{Source and navigation}\label{source-and-navigation}

\href{Continuous\%20Notes/Computing\%20Foundations\%20-\%20From\%20Bits\%20to\%20C++\%20and\%20Systems\%20APIs\#From\%20three\%20repairs\%20to\%20five\%20decisions\%20to\%20zero\%20boilerplate}{Teaching mechanism and examples} · \hyperref[appendix-fd090d982a15]{Exact coverage} · \href{}{Dependency Map (02-Foundations-II, appendix-4ad47bdd269a)}
